\documentclass[12pt]{article}
\ProvidesPackage{preamble}

\renewcommand{\familydefault}{\sfdefault}
\usepackage{lmodern}

% Language setting
% Replace `english' with e.g. `spanish' to change the document language
\usepackage[english]{babel}

% Set page size and margins
% Replace `letterpaper' with `a4paper' for UK/EU standard size
\usepackage[letterpaper,top=1in,bottom=1in,left=1in,right=1in,marginparwidth=1.75cm]{geometry}

% Useful packages
\usepackage{amsmath,amsthm,amssymb,color,latexsym}
\usepackage{mathtools}
\usepackage{braket}
\usepackage{graphicx}
\usepackage{esint} % for the wacky integrals
\usepackage[colorlinks=true, allcolors=blue]{hyperref}
\usepackage{subfig}
\usepackage{mathtools}
\usepackage{booktabs}
\usepackage[all]{xy}

\usepackage[shortlabels]{enumitem}
\usepackage[table]{xcolor}
\usepackage{geometry}
\usepackage{url}
\geometry{letterpaper}
\usepackage{parskip}

\setlength{\parindent}{0pt}

\newtheorem{problem}{Problem}

\newtheorem{theorem}{Theorem}
\newtheorem*{proposition}{Proposition}
\newtheorem{lemma}[theorem]{Lemma}
\newtheorem{corollary}[theorem]{Corollary}
\newtheorem{conjecture}[theorem]{Conjecture}
\newtheorem{postulate}[theorem]{Postulate}

\theoremstyle{definition}
\newtheorem{definition}[theorem]{Definition}
\newtheorem{ex}[theorem]{Example}

\theoremstyle{remark}
\newtheorem*{remark}{Remark}
\newtheorem*{notation}{Notation}
\newtheorem*{note}{Note}

\newenvironment{solution}[1][Solution]{\textbf{\textit{#1. }}}{$\blacksquare$}

% shortcuts
\newcommand{\ndiv}{\not\hspace{2.5pt}\mid}
\newcommand{\NN}{\mathbb{N}}
\newcommand{\ZZ}{\mathbb{Z}}
\newcommand{\QQ}{\mathbb{Q}}
\newcommand{\RR}{\mathbb{R}}
\newcommand{\CC}{\mathbb{C}}

\newcommand{\sS}{\mathcal{S}}
\newcommand{\sT}{\mathcal{T}}
\newcommand{\ps}{\mathcal{P}}

\begin{document}
\noindent MAT240 09-17 \hfill Sky Hong

\hrule

%%%%%%%%%%%%%%%%%%%%%%%%%%%%%%%%%%%%%%%%%%%%%%%%%%%%%%%%%%%%%

Go to tutorial!!! HW1 is due sunday at 8pm. Mathematics SUpport Initiative, see Quercus, math help club.

Recap. Recall that $\ZZ / m \ZZ = \{ \bar{0}, \dots, \overline{m-1} \}$is the set of equivalence classes mod $m$, equipped with the operations $+$ and $\cdot$, given by:
\begin{enumerate}
    \item Do the operation as usual
    \item Then, take the remainder
\end{enumerate}

\begin{remark}
    What is a prime number?
\end{remark}

\begin{defn}
    Let $p \in \NN$, $p > 1$, then $p$ is prime if $\forall x,y \in \ZZ$,\[
        p \nmid x \text{ and } p \nmid y \implies p \nmid xy
    \]
\end{defn}

asfasf
\begin{remark}
    What is a prime number?
\end{remark}

Check this as an exercise! This is also called Euclid's lemma.

\begin{thm}
    Let $m \in \NN$. If $m$ is prime, then $\ZZ / m \ZZ$ is a field. blahblahblahblahblahblahblah blah blah blah blah
\end{thm}
\begin{remark}
    What is a prime number?
\end{remark}
\begin{pf}
    First, check the field axioms, with (F4) holding for addition.
\end{pf}

\end{document}